\section{ResNet：残差学习}

\subsection{深度退化问题}

\begin{figure}[H]
\centering
\includegraphics[width=0.95\textwidth]{slides-images/slide-027.jpg}
\caption{深度退化现象：56 层普通网络的训练和测试误差都高于 20 层网络\protect\footnotemark}
\end{figure}
\footnotetext{Slides 第28页。}

一个令人困惑的发现：当把普通的 CNN（没有残差连接）从 20 层加深到 56 层时，\textbf{训练误差和测试误差都变高了}。

\begin{warningbox}{这不是过拟合！}
如果是过拟合，我们会期望看到：训练误差下降（模型``记住''了训练数据），但测试误差上升。然而实际观察到的是\textbf{训练误差也更高}——56 层网络在训练集上的表现甚至不如 20 层网络。

这说明问题出在\textbf{优化}上，而非模型容量上。更深的网络包含了更浅网络能学到的所有函数（理论上可以将多余的层设为恒等映射），但优化算法\textbf{找不到}好的解。
\end{warningbox}

\subsection{残差连接的核心思想}

\begin{figure}[H]
\centering
\includegraphics[width=0.95\textwidth]{slides-images/slide-029.jpg}
\caption{残差块：$H(x) = F(x) + x$，其中 $F(x)$ 是需要学习的残差映射\protect\footnotemark}
\end{figure}
\footnotetext{Slides 第30页。}

ResNet 的核心创新：不让网络直接学习目标映射 $H(x)$，而是学习\textbf{残差} $F(x) = H(x) - x$。网络输出为：

$$H(x) = F(x) + x$$

\begin{importantbox}{残差连接解决了什么问题？}
\begin{enumerate}
    \item \textbf{恒等映射变得容易}：如果某层不需要做任何变换，网络只需让 $F(x) \approx 0$（将滤波器权重学到接近零），输出就自然等于输入 $x$。
    \item \textbf{学习增量而非全量}：网络从``从头学习完整映射''变成``在输入基础上学习一个微小的修正量''，这个优化问题更容易。
    \item \textbf{梯度传播更顺畅}：反向传播时，梯度可以通过跳跃连接直接回传到前面的层，缓解了梯度消失问题。
\end{enumerate}
\end{importantbox}

\subsection{残差块的结构}

一个标准的残差块包含：
\begin{enumerate}
    \item 输入 $x$
    \item $3 \times 3$ 卷积 $\rightarrow$ ReLU $\rightarrow$ $3 \times 3$ 卷积 得到 $F(x)$
    \item 跳跃连接：将 $x$ 直接复制过来
    \item 相加：$F(x) + x$
    \item ReLU
\end{enumerate}

\begin{figure}[H]
\centering
\includegraphics[width=0.95\textwidth]{slides-images/slide-031.jpg}
\caption{ResNet 的完整架构：由多个残差块堆叠而成，周期性地下采样并翻倍通道数\protect\footnotemark}
\end{figure}
\footnotetext{Slides 第32页。}

\begin{warningbox}{残差连接要求张量形状一致}
$F(x) + x$ 要求两者形状完全相同。这就是为什么残差块内部使用 $3 \times 3$ 卷积、stride 1、padding 1——保持空间尺寸不变。当需要下采样或改变通道数时，需要在跳跃连接上也做相应的变换（通常用 $1 \times 1$ 卷积）。
\end{warningbox}

\subsection{ResNet 的设计模式}

\begin{figure}[H]
\centering
\includegraphics[width=0.95\textwidth]{slides-images/slide-033.jpg}
\caption{ResNet 家族：18 层到 152 层，性能随深度提升，但收益递减\protect\footnotemark}
\end{figure}
\footnotetext{Slides 第34页。}

ResNet 的整体设计：
\begin{itemize}
    \item 入口：一个较大的 $7 \times 7$ 卷积层（stride 2）+ max pooling
    \item 主体：多个\textbf{阶段}（stage），每个阶段包含若干残差块
    \item 阶段之间：空间尺寸减半，通道数翻倍
    \item 出口：全局平均池化 $\rightarrow$ 全连接层 $\rightarrow$ 分类
\end{itemize}

ResNet 家族提供了不同深度的变体：
\begin{table}[H]
\centering
\begin{tabular}{lcc}
\toprule
\textbf{模型} & \textbf{层数} & \textbf{Top-5 错误率} \\
\midrule
ResNet-18 & 18 & 10.9\% \\
ResNet-34 & 34 & 8.6\% \\
ResNet-50 & 50 & 7.1\% \\
ResNet-101 & 101 & 6.4\% \\
ResNet-152 & 152 & 6.2\% \\
\bottomrule
\end{tabular}
\caption{ResNet 家族在 ImageNet 上的性能}
\end{table}

\begin{importantbox}{ResNet 的历史影响}
ResNet（2015）是第一个成功训练超过 100 层的深度网络。它赢得了 ILSVRC 2015 的冠军，首次在 ImageNet 上超越了人类水平。更重要的是，残差连接不仅用于 CNN——\textbf{现代 Transformer 架构中也广泛使用残差连接}，原因完全相同：帮助梯度在非常深的网络中有效传播。
\end{importantbox}

\begin{knowledgebox}{Kaiming He 的贡献}
ResNet 的第一作者何恺明（Kaiming He）是计算机视觉领域最有影响力的研究者之一。除了 ResNet 之外，他还发明了 Kaiming 初始化方法（本节课后面会讲到），两者同出一人之手——这并非巧合，因为残差连接和正确的初始化都是为了解决\textbf{深度网络的训练稳定性}问题。
\end{knowledgebox}

\subsection{本章小结}

ResNet 通过残差连接解决了深层网络的优化困难。残差学习的核心思想是``学习增量而非全量''——让网络在输入的基础上学习一个小的修正量。这个看似简单的改变使得训练超过 100 层的网络成为可能，并且其影响远超 CNN 领域，成为现代深度学习架构的标准组件。

%% ========== Part 7: 权重初始化 ==========

